% element: 10-s008:e04
\BeleskeID{10-s008}
\subsection*{Uključivanje}
U intervalu $t_0$–$t_1$ pune se oblasti prostornog tovara koje su bile inverzno polarizovane; njihove širine se smanjuju. Spoj je još inverzno polarizovan, nema tranzistorskog efekta i kolektorska struja je nula.

U $t_1$–$t_2$ nastavlja se punjenje ekvivalentnih kapacitivnosti. Bazno-emitorski spoj postaje direktno polarizovan, oblast je uska i u bazi se pojavljuje višak manjinskih nosilaca. Postoji tranzistorski efekat: kolektorska struja zavisi od bazne, a izlazni napon opada.

Pri kraju tog intervala oba spoja, baza–emitor i baza–kolektor, direktno su polarizovana. U $t_2$–$t_3$ postoji veliki višak manjinskih nosilaca na oba spoja; tranzistor je u zasićenju.
\subsection*{Isključivanje}
U $t_3$–$t_4$ ulazni napon je negativan, ali viškovi nosilaca u oblasti baze ne mogu trenutno da nestanu. Negativna bazna struja ih eliminiše, dok tranzistor još ostaje u zasićenju.

U $t_4$–$t_5$ spoj baza–kolektor postaje inverzno polarizovan i oblast nastavlja da se širi. Opada broj nosilaca uz spoj baza–emitor, kao i kolektorska struja.

U $t_5$–$t_6$ spoj baza–kolektor je inverzno polarizovan; i spoj baza–emitor postaje inverzno polarizovan i širi se. Punjenje oblasti prostornog tovara, odnosno njihovih kapacitivnosti, objašnjava eksponencijalan oporavak.
